% =============================================================================
% 8. CONCLUSION
% =============================================================================
\section{Conclusion}
\label{sec:conclusion}

We started from the observation that a single step of a coding agent and
a single function call site share the same four-part conditioning
structure---\textit{context, action, externally produced return,
continuation}---and that this structural isomorphism makes ordinary
source code a natural, internet-scale supply of agent-relevant signal.
Building on this framing, we introduced function-aware FIM mid-training:
a self-supervised stage that masks targets selected through program
dependency graph analysis and a complexity--inferability double
criterion, with chain-of-thought rationales embedded inside the FIM
middle span. Across three independent axes of robustness---model size
(7B/14B/32B), post-training pipeline (r2e-gym, SWE-Smith, swe-lego), and
base-model family (Qwen2.5-Coder, Qwen3)---the recipe delivers
consistent gains on coding-agent benchmarks, and the same Python-only
mid-training corpus also transfers to non-coding tool-use benchmarks
($\tau$-bench, BFCL) where it has no a priori reason to help, providing
direct evidence for the motivating isomorphism. Promising directions for
future work include extending the function-aware selection pipeline to
non-Python languages, validating the mid-training stage on additional
base-model families beyond the Qwen lineage, and studying how
mid-training composes with reinforcement-learning-based post-training,
where a structural prior over function-call-shaped trajectories may
shape the exploration landscape in ways that supervised post-training
alone cannot.
